\clearpage
\section{Contraction idempotence in all observing contexts}
\label{sec:copying-contraction-idempotence}
The next completed structural dependency of Proposition~12.2 is the Idem
family of Figure~21, for both term and atom sharing. Two weakenings feeding
a contraction are contextually equivalent to one weakening of the same sort.
The proof establishes output closure, input safety, full forward and converse
robustness, and exact successful-transition counts in every original closing
context. Observing sites may be nested under arbitrarily many boxes. No
binding-free restriction is imposed on the observation class.

\begin{theorem}[Contraction-idempotence contextual equivalence]
Let $s\in\{\mathsf{term},\mathsf{atom}\}$, let $W_s$ be weakening and $C_s$
contraction, and set
\[
 L_s=(W_s\otimes W_s);C_s : \epsilon\Rightarrow s,
 \qquad R_s=W_s : \epsilon\Rightarrow s.
\]
At every original addressed occurrence in a focus-free observing graph,
if the two filled graphs are closed, replacing $L_s$ with $R_s$ preserves
successful termination at every exact number of machine transitions.
The same holds for arbitrary simultaneous recursive Idem fillings.
\end{theorem}
The original graph interface and the two literal fragments are retained.
The theorem is applied to actual fragment replacement, since these fragments
have no inputs and are not themselves one-term-input open programs.
In particular, it does not follow from a representation identification:
the left graph has three edges and the right graph one.

\paragraph{Why this milestone advances the example.}
Table~3 makes Idem a prerequisite of pulling weakening out of a box
($\mathrm{BPullW}$), while Section~12.3 uses these two laws to prove semantic
weakening. Those weakening laws extend the closed Figure~22 example to the
arbitrary distinct-list environments in Proposition~12.2. Idem also supports
box-contraction pulling and parametricity. This section completes its full
template obligation; box-pulling, semantic weakening and the remaining
Figure~22 links are still separate tasks.

\paragraph{Correspondence with the paper's proof.}
Appendix~H.3.1 proves Idem robustness by choosing a maximal applicable sharing
tree and applying the tree-preservation transfer property of Corollary~H.8(3)
(HTML Corollary~47(3)). Here the same robustness strategy is implemented using
proved successor coherence and the already constructed selected-path
occurrences. Rematching a smaller occurrence leaves uninspected sibling
branches attached at its boundary; coherence identifies its result with the
independently supplied actual successor. No copying or compute rule changes.
The nullary argument below explains why this smaller occurrence avoids every
shallow Idem replacement. This is a difference in proof organization, not in
the equation or its observation class.

\subsection{The literal contraction-idempotence template}
\termdef{literal-contraction-idempotence}
Fix a sharing sort $s$, either term or atom. Write $w_s$ for the actual
weakening corolla with interface $[]\to[s]$, and $c_s$ for the actual
contraction corolla with interface $[s,s]\to[s]$. The two graphs of
Figure~21's Idem family are
\[
 L_s=(w_s\otimes w_s);c_s,
 \qquad R_s=w_s.
\]
Here tensor is ordered juxtaposition and the semicolon is the existing
serial graph constructor. Thus $L_s$ contains two distinct weakening edges,
the two wires joining them to the ordered contraction inputs, and the
contraction itself. Both graphs have empty input interface and one output
of sort $s$. These definitions retain the literal graphs at every depth
allowance; no distributor or translated term is simplified in their definition.

\lean{Copying/ContractionIdempotence}{14}{16}{ContractionIdempotence.left}
\lean{Copying/ContractionIdempotence}{19}{20}{ContractionIdempotence.right}
\par\penalty0


The relation called \texttt{Step} contains these two sort-indexed instances.
Despite its name, it is a structural graph-pair relation, not a machine
transition. Its closure \texttt{UpToIso} relates $G$ to $H$ when $G$ is
recursively isomorphic to some literal left side and the corresponding
literal right side is recursively isomorphic to $H$. The two carrier
renamings are independent and preserve every ordered interface position.
Neither relation assumes the desired contextual equivalence.

\lean{Copying/ContractionIdempotence}{23}{25}{ContractionIdempotence.Step}
\lean{Copying/ContractionIdempotence}{29}{32}{ContractionIdempotence.UpToIso}
\par\penalty0


If a predicate on primitive labels holds of weakening and contraction of
sort $s$, it holds throughout $L_s$: first combine the two weakening
corollas by the tensor label lemma, then append the contraction by the
serial label lemma. On $R_s$ the weakening hypothesis alone suffices.
Specializing these facts to exclusion of focus and hole primitives proves
both graphs focus-free and hole-free. Inspection of the single rule
constructor proves its input interface is empty. These three properties
pass to \texttt{UpToIso}: recursive isomorphisms preserve label predicates,
and the external interfaces are already fixed by the isomorphism types.

\lean{Copying/ContractionIdempotence}{36}{43}{ContractionIdempotence.left_allLabels}
\lean{Copying/ContractionIdempotence}{46}{48}{ContractionIdempotence.right_allLabels}
\lean{Copying/ContractionIdempotence}{50}{55}{ContractionIdempotence.Step.focusFree}
\lean{Copying/ContractionIdempotence}{57}{62}{ContractionIdempotence.Step.holeFree}
\lean{Copying/ContractionIdempotence}{65}{67}{ContractionIdempotence.Step.input}
\lean{Copying/ContractionIdempotence}{69}{73}{ContractionIdempotence.UpToIso.focusFree}
\lean{Copying/ContractionIdempotence}{75}{79}{ContractionIdempotence.UpToIso.holeFree}
\lean{Copying/ContractionIdempotence}{81}{83}{ContractionIdempotence.UpToIso.input}
\par\penalty0


The Idem pre-template uses precisely \texttt{UpToIso} as its relation and
the proved focus-freedom statement as its structural certificate. Taking
identity isomorphisms exhibits every literal pair as a template instance.

\lean{Copying/ContractionIdempotence}{87}{89}{ContractionIdempotence.template}
\lean{Copying/ContractionIdempotence}{92}{93}{ContractionIdempotence.literal}
\par\penalty0


\paragraph{Admissibility and representation allowance.}
The empty-interface theorem proves \texttt{NoInputs}, and hence the weaker
\texttt{NoTermInputs} property. Output closure follows because an
operation-only path starting at an external term input cannot exist.
Input safety is equally direct: an entering token would supply a source
port of a template hole, contradicting that hole's empty input interface.
This argument works for every allowed-context class and every pair of
cost orders; it assumes no simulation or robustness theorem.

\lean{Copying/ContractionIdempotenceLaws}{11}{11}{ContractionIdempotence.template_noInputs}
\lean{Copying/ContractionIdempotenceLaws}{13}{13}{ContractionIdempotence.template_noTermInputs}
\lean{Copying/ContractionIdempotenceLaws}{15}{15}{ContractionIdempotence.template_outputClosed}
\lean{Copying/ContractionIdempotenceLaws}{19}{22}{ContractionIdempotence.template_inputSafe}
\par\penalty0


An extra representation allowance preserves the single weakening literally.
For the left side, commute padding through serial assembly, tensor assembly,
and each primitive corolla. The serial constructor may use a different
internal carrier presentation, so the result is an ordered recursive
isomorphism rather than an asserted literal graph equality. These two
results show that a padded literal pair is again an Idem pair up to
isomorphism. Compose with the padded isomorphisms of an arbitrary existing
pair to prove \texttt{UpToIso.raise}, and therefore padding closure of the
pre-template. All three left-side edges survive this construction.

\lean{Copying/ContractionIdempotenceLaws}{26}{31}{ContractionIdempotence.left_raise}
\lean{Copying/ContractionIdempotenceLaws}{34}{35}{ContractionIdempotence.right_raise}
\lean{Copying/ContractionIdempotenceLaws}{39}{44}{ContractionIdempotence.Step.raise}
\lean{Copying/ContractionIdempotenceLaws}{46}{52}{ContractionIdempotence.UpToIso.raise}
\lean{Copying/ContractionIdempotenceLaws}{54}{54}{ContractionIdempotence.template_paddingClosed}
\par\penalty0



\subsection{Nullary pre-templates and their safety}
\termdef{nullary-template}
A pre-template $T$ is \emph{nullary} when every related pair $G\,T\,H$ has
literally the empty common input interface. This is a property of its graph
relation. Exchanging endpoints preserves it. It implies the previously
formalized no-term-input property, since an index into the empty input list
cannot exist. Output closure follows: there is no operation path from a term
input to an output. The original input-safety condition is consequently
vacuous for any context class and any cost parameters. These conclusions
supply their proofs to the existing specimen machinery; they do not assume
execution equivalence.
\lean{Copying/NullaryTemplates}{10}{11}{PreTemplate.NoInputs}
\lean{Copying/NullaryTemplates}{14}{15}{PreTemplate.NoInputs.converse}
\lean{Copying/NullaryTemplates}{18}{23}{PreTemplate.NoInputs.noTermInputs}
\lean{Copying/NullaryTemplates}{26}{27}{PreTemplate.NoInputs.outputClosed}
\par\penalty0
\lean{Copying/NullaryTemplates}{31}{34}{PreTemplate.NoInputs.inputSafe}
\subsection{A nullary replacement cannot absorb a sourced footprint}
\termdef{nullary-ancestry-retention}
Fix genuine recursive filling certificates for two expansions $C[L]$ and
$C[R]$ of a nullary pre-template. Suppose an actual expanded edge consumes
a vertex which is the image of an original vertex of $C$. That edge cannot
belong to a shallow replacement: such an incidence would supply an input
index of the replacement, contradicting its empty input interface. It must
therefore come from a retained outer edge. Its target vertices are images
of original vertices as well.

Induction along an actual edge ancestry repeats this argument at each next
source incidence. It proves that every encountered outer edge is retained,
and reflects the entire ancestry into the original context. When the left
expanded graph is hole-free, a retained original hole cannot remain
unfilled. A source at such a hole would again supply an impossible input.
Thus neither the original sourced edge nor any of its descendants lies
in a shallow hole. The proof does not require a rooted history, binding
freedom, or an acyclic surrounding context.
\lean{Copying/NullaryFootprint}{14}{20}{NullaryOverlap.source_retained}
\lean{Copying/NullaryFootprint}{24}{36}{NullaryOverlap.ancestor_retained}
\lean{Copying/NullaryFootprint}{39}{45}{NullaryOverlap.source_not_hole}
\lean{Copying/NullaryFootprint}{49}{60}{NullaryOverlap.ancestor_not_hole}
\par\penalty0

\subsection{Selected atom trees for all store rules}
Read follows an instance edge through the selected atom-sharing path to its
store. Its reader consumes the original token vertex, so the preceding
retention theorem applies to the reader and every edge downstream from it.
A selected subtree occurrence carries an actual ancestry from its inspected
input to each selected sharing edge. Extending it through the unchanged read
frame therefore retains the whole rematched source, including its instance
and store, while leaving irrelevant sibling branches outside that source.

Assignment has the same selected store path and also the original stable
new value. Every edge of that value consumes the value's sole term input,
so all of it is directly downstream of the assignment head. Its complete
ordered output interface, including thunk outputs, is retained. This proves
the two selected-source exclusion lemmas.
\lean{Copying/NullaryStoreFootprint}{20}{50}{NullaryOverlap.read_subtree_avoids}
\lean{Copying/NullaryStoreFootprint}{54}{92}{NullaryOverlap.assignment_subtree_avoids}
For equality of two references resolving to one store, the two selected
inputs may merge at a sharing contraction. The genuine two-input subtree
contains the union of their ancestries. Both instance edges are downstream
of the comparator; every selected sharing edge is downstream of at least
one instance, and the store is downstream of a successful resolution.
For references to different stores, make the same argument separately in
the two atom trees. The complete occurrence preserves distinct store edges,
the false decision, and the order of the two content outputs. These are
retainedness proofs for the original machine sources, not new store rules.
\lean{Copying/NullaryStoreFootprint}{96}{149}{NullaryOverlap.same_subtree_avoids}
\lean{Copying/NullaryStoreFootprint}{155}{221}{NullaryOverlap.different_subtrees_avoids}
Finally, discharge all selected-subtree premises from arbitrary original
rule occurrences. For read and assignment take the minimal spine from input
zero to the root, extend its full-label occurrence through the corresponding
frame, and compose with the original outer match. For same-store equality
use two-input subtree extraction with the original inputs zero and one
first in the new boundary order. For different-store equality take one
minimal spine in each tree. In every case the active comparison or operation
input maps to the original token vertex, and the preceding lemmas prove
avoidance of shallow holes. No smaller match is supplied by the caller.
\lean{Copying/NullaryStoreFootprint}{227}{240}{NullaryOverlap.read_occurrence}
\lean{Copying/NullaryStoreFootprint}{244}{260}{NullaryOverlap.assignment_occurrence}
\lean{Copying/NullaryStoreFootprint}{264}{274}{NullaryOverlap.same_occurrence}
\lean{Copying/NullaryStoreFootprint}{278}{294}{NullaryOverlap.different_occurrence}
\par\penalty0

\subsection{Pointed residuals in unrestricted contexts}
\termdef{unrestricted-pointed-residual}
The finite-cut construction already supplies a genuine recursive residual
witness with the computed graft as its common-context geometry. A local
common-context resume vertex has an explicit image in this graft. Applying
the original unrestricted quasi-specimen construction at that vertex gives
the two actual resumed states. All matched and untouched hole names, nested
box bodies and ordered ports are preserved by the established pointed
correspondence lemmas. Since the allowed-context predicate is true, no
binding-path condition is needed on this graft.

The unrestricted input diamond first lifts the actual local rewrite to the
left graph and the constructed paired local rewrite to the right graph.
Both return through the same local input coordinate. Machine preservation
proves hole freedom of both resulting whole graphs; the pointed residual
then gives the original quasi-specimen at their exact search tokens.
This is the common transfer interface for all nine cases below.
\lean{Copying/AllContextResiduals}{10}{29}{RetainedFootprint.pointedResidual_allContexts_quasi}
\lean{Copying/AllContextResiduals}{34}{58}{RetainedFootprint.allContexts_input_diamond}
\subsection{Primitive footprints and unrestricted residuals}
Let $T$ be a pre-template, let $C$ be an original common context with
recursive filling families $L,R$, and suppose the left expansion is
hole-free. In the footprint results below, $T$ has no inputs. A complete
occurrence $f:A\to C[L]$ preserves all primitive and recursive labels,
incidence rows and edge identities. Suppose its selected local source
port attaches to an original context vertex, and every edge of $A$ is a
descendant of an edge consuming that port. The nullary ancestry barrier
then excludes every such edge from an original shallow hole. Map each
local ancestry path through $f$, starting with the mapped source incidence,
and apply that barrier. This proves avoidance for the complete footprint;
no binding freedom, rooted execution history or non-entry assumption is
supplied to this lemma.

\lean{Copying/NullaryPrimitiveOverlap}{18}{28}{NullaryOverlap.sourced_ancestors_avoid}
\par\penalty0


For allocation, the source consists of the reference head and the original
stable argument. The head is its own ancestor. Each argument edge consumes
the stable argument's input, which is attached directly to the reference
head's output, so it is a one-edge descendant. For unary arithmetic, the
only other edge is the numeric operand; the existing exact adjacency
classification supplies its path. For binary arithmetic, the same
classification covers both ordered numeric operands and either original
binary arithmetic code. Applying the preceding general result gives
complete retained footprints for all three cases, retaining the literal
stable operand and operation parameters.

\lean{Copying/NullaryPrimitiveOverlap}{32}{50}{NullaryOverlap.allocation_avoids}
\lean{Copying/NullaryPrimitiveOverlap}{53}{61}{NullaryOverlap.unary_avoids}
\lean{Copying/NullaryPrimitiveOverlap}{65}{74}{NullaryOverlap.binary_avoids}
\par\penalty0


\termdef{nullary-primitive-residual}
The residual theorem is slightly more general than these nullary footprint
lemmas. It accepts any certified retained complete primitive footprint,
not a nullarity hypothesis. Let the original local rewrite be
\[
 A,\operatorname{input}(i)\longrightarrow
 B,\operatorname{input}(i),
\]
where input $i$ has term type. Primitive rigidity identifies the actual
right-filled local source with $A$ up to recursive isomorphism, so it
constructs the same original local rewrite on the right. It does not
assume that a paired local rewrite exists. Focus freedom of $C$ and the
recursive certificates imply focus freedom of $C[L]$; restricting through
the complete occurrence proves $A$ focus-free, and actual rewrite
preservation then proves $B$ focus-free.

Use the reflexive filling witness on $B$ as the common local residual.
The existing complete cut and pointed graft constructions give actual
rewrites of $C[L]$ and $C[R]$, preserving every unmatched original filling.
They resume at the same original local input embedded into their respective
grafts. Their whole successors are hole-free by rewrite preservation.
The common residual is admissible in the unrestricted context class, so
no guarded-output condition or binding-free cut hypothesis is required.
The conclusion contains both actual whole-graph rewrites and an
identity-refined quasi-specimen for their resumed search states.

\lean{Copying/NullaryPrimitiveResiduals}{23}{42}{RetainedFootprint.primitive_allContexts_input_diamond}
\par\penalty0


Finally, compare an arbitrary supplied whole-machine successor with this
constructed left rewrite using the already proved successor determinism.
The left needs zero additional transitions and the right takes its
constructed single transition. Thus the local cost requirement is
$Q(1+0,1)$, available for every cost order by reflexivity. The quasi-specimen
can be enlarged from identity refinements to the two requested execution
refinements, again by reflexivity. This proves the full quantitative case,
not merely an existence statement for a conveniently chosen successor.
Exchanging the two actual rewrites proves the corresponding converse case
in the opposite of the original specimen system. The exchange uses no
symmetry assumption on $T$.

\lean{Copying/NullaryPrimitiveResiduals}{46}{62}{RetainedFootprint.primitive_allContexts_robust_case}
\lean{Copying/NullaryPrimitiveResiduals}{66}{83}{RetainedFootprint.primitive_allContexts_converse_robust_case}
\par\penalty0
\subsection{Automatic binding freedom for nullary fillings}
\termdef{nullary-context-binding-freedom}
Let $T$ be a pre-template whose related graphs have empty input interfaces.
Although the observation class is unrestricted, a common context that gives
a closed recursive $T$-filling has a useful derived property: it is
binding-free. This statement is about the actual certified context, and
imposes no additional restriction on the class of observations.

First, any positive edge-ancestry path ending at an edge provides a source
of that final edge. Prove this by induction along the path. If the final
edge is a shallow hole of the common context, hole freedom of the expanded
left graph excludes an unfilled retained hole. Its certificate therefore
contains an actual $T$-pair. Nullarity makes its input interface empty,
so the common context's typing law makes its source list empty. This
contradicts the path's source witness and proves shallow binding freedom.
The same argument shows that no common-context token can enter such a hole,
regardless of the token's mode or wire type.

\lean{Copying/NullaryBindingFree}{13}{22}{NullaryOverlap.strictAncestor_source}
\lean{Copying/NullaryBindingFree}{26}{36}{NullaryOverlap.shallow_bindingFree}
\lean{Copying/NullaryBindingFree}{40}{49}{NullaryOverlap.not_entering}
For recursive binding freedom, induct on the depth allowance. At depth
zero the preceding shallow argument suffices. At positive depth, retained
box corollas have hole-free bodies because the actual expanded left graph
is hole-free. Filled shallow holes cannot be boxes. For an originally boxed
recursive certificate, the outer focus-freedom assumption and the actual
closed filled body give exactly the induction hypotheses. Applying this
result through a filling witness's left endpoint isomorphism establishes
binding freedom of its literal common context.

\lean{Copying/NullaryBindingFree}{54}{83}{DeepFilling.nullary_bindingFree}
\lean{Copying/NullaryBindingFree}{87}{89}{FillingWitness.nullary_bindingFree}
\subsection{Actual beta and copying footprints}
Every outer beta edge has a forward path from its application head. The
lambda is adjacent to the application, the box is adjacent to the lambda,
and each edge in the stable argument has the argument's sole term input
as a source. Mapping these paths through the original complete occurrence
and applying nullary source-ancestry retention proves that the entire beta
footprint avoids shallow replacements. The body inside the box remains
arbitrary and may contain recursive template occurrences.

For copying, derive binding freedom and non-entry from the nullary filling,
then apply the already proved minimal-spine construction. Its result is an
actual occurrence of the original copying rule with a selected contraction
spine and the complete copyable payload. All selected edges are retained;
uninspected sibling branches of the originally supplied tree stay in the
computed complement. Normalizing the payload uses only a proved recursive
representation isomorphism. It changes neither the copying rule nor the
surrounding graph.

\lean{Copying/NullaryDeepOverlap}{13}{52}{NullaryOverlap.beta_avoids}
\lean{Copying/NullaryDeepOverlap}{56}{69}{NullaryOverlap.minimal_copy_occurrence}
\subsection{Local reducts and unrestricted residual specimens}
The preceding automatic binding-freedom result permits reuse of the generic
retained-decoration construction. In the beta case, extract the actual
paired box body, pad it using the template's padding-closure theorem, expose
its bound output, and splice the unchanged stable argument. In the copying
case, inspect every original copyability constructor; primitive payloads
are rigid and lambda payloads duplicate their actual recursively paired
bodies. These constructions supply genuine local rewrites in the paired
source expansion and filling witnesses between their reducts. The exact
host-port equation identifies the resumed input on both sides.

\lean{Copying/NullaryDeepLocal}{12}{27}{NullaryResidual.beta_local}
\lean{Copying/NullaryDeepLocal}{31}{46}{NullaryResidual.copying_local}
The final beta and copying cases place these local diagrams in the computed
common complement. The unrestricted input diamond constructs both actual
whole-graph rewrites and an equality-up-to quasi-specimen for the original
all-context system. No graft-admissibility premise is needed in this context
class. The copying case first constructs its retained smaller occurrence;
no footprint certificate is supplied by the caller.

Each quantitative theorem accepts an independently supplied actual
representation-class successor. Proved successor determinism identifies it
with the constructed left reduct. The comparison uses zero additional left
steps and one actual right rewrite, followed by the constructed residual.
Thus these branches work for arbitrary cost preorders; taking exact cost
later yields the required count-preserving robustness. There are no
rootedness, binding-freedom or overlap assumptions in these public branch
theorems. Beta requires template padding closure, and copying keeps the
original copyability premise without strengthening it.

\lean{Copying/NullaryDeepResiduals}{12}{32}{NullaryResidual.beta}
\lean{Copying/NullaryDeepResiduals}{36}{58}{NullaryResidual.copying}
\subsection{Universal robustness in both directions}
For any nullary padding-closed pre-template, the actual local-rule case
analysis now covers all nine original constructors. Allocation and arithmetic
use their proved complete retained footprints. Read, assignment and the two
equality decisions first construct the selected atom-tree occurrences.
Deep beta and copying use the actual recursive decoration and paired local
residual just proved. Original stability and copyability hypotheses pass
unchanged to the original rules. Each case relates any independently
supplied source successor to a constructed opposite-side rewrite using zero
extra source transitions and one target transition.

An arbitrary contextual rewrite yields its own complete local occurrence,
which discharges the extraction premises of that case analysis. An arbitrary
representation-class successor reflects to a finite actual step. Search is
impossible in rewrite mode; hence that step is one of the covered rewrites.
\lean{Copying/NullaryRobustness}{21}{69}{NullaryResidual.local_case}
\lean{Copying/NullaryRobustness}{72}{83}{NullaryResidual.rewrite_case}
\lean{Copying/NullaryRobustness}{87}{100}{NullaryResidual.class_step}
The forward universal theorem applies this result to the original specimen's
context, recursive certificates, focus freedom and closed expanded graphs.
The existence of a matching transition and residual is proved for every
successor, with no extra match, ancestry, non-entry or binding-freedom premise.

For the converse, reverse the actual filling certificates. Nullarity and
padding closure hold in the reversed direction. The same exhaustive argument
constructs an ordinary forward transition on the original left side.
Convert its reversed-template quasi-specimen to the opposite of the original
specimen system by swapping filling families in the unchanged context. The
conversion preserves arbitrary refinement-up-to relations in their original
orientation. Both universal conclusions hold for arbitrary cost orders;
their substantive local comparison is reflexivity at $Q(1+0,1)$.
\lean{Copying/NullaryRobustness}{105}{111}{Specimens.converse_quasi_rel}
\lean{Copying/NullaryRobustness}{115}{137}{PreTemplate.NoInputs.robust}
\lean{Copying/NullaryRobustness}{140}{165}{PreTemplate.NoInputs.robust_converse}
\subsection{Exact execution in unrestricted addressed contexts}
A closed recursively filled pair supplies an initial lifting for the original
all-context specimen system. Construct its specimen at the common external
term input in search mode. The filling witness's actual endpoint isomorphisms
identify the two filled states with the stated graphs, and both are rooted
by the empty initial search. This initial-lifting result is generic and does
not assume nullarity or binding freedom.

\lean{Copying/NullaryExecution}{12}{27}{FillingWitness.allContexts_initial_lifting}
For a padding-closed nullary template, its proved input safety and full
forward and converse robustness discharge the hypotheses of sufficiency
with the exact cost triple. Therefore any original closed paired filling
has exact quantitative execution equivalence. Unpack the two refinements:
if one side has a successful run of length $n$, the other refinement supplies
a successful run whose length must also be $n$. Applying the reverse
refinement proves the converse. Neither successful termination nor any
behaviour of the surrounding computation is a premise.

\lean{Copying/NullaryExecution}{35}{43}{NullaryExecution.filling_exact_equiv}
\lean{Copying/NullaryExecution}{47}{60}{NullaryExecution.filling_termination_count}
To obtain an observation theorem for an individual nullary fragment, take
any original addressed site in a focus-free observing graph with external
interface $\star\Rightarrow\epsilon$. The site's depth, input interface and
output interface are arbitrary. A genuine template pair supplies a filling
witness at the site, and addressed closure rebuilds each enclosing box and
retains every other context edge. If the two resulting filled graphs are
hole-free, the preceding exact theorem applies. This quantifies over all
such original contexts, including sites beneath arbitrarily many boxes,
and imposes no binding-free restriction. The addressed count theorem states
the same conclusion separately for every successful transition count.

\lean{Copying/NullaryExecution}{65}{70}{NullaryExecution.addressed_exact_equiv}
\lean{Copying/NullaryExecution}{74}{79}{NullaryExecution.addressed_termination_count}
\subsection{The complete contraction-idempotence theorem}
\termdef{contraction-idempotence-equivalence}
Instantiate the generic results with the literal Figure~21 Idem template.
Its proved empty input interface discharges nullarity, and its proved
padding law discharges the only additional closure hypothesis. Consequently
Idem and its converse are robust for all original unrestricted specimens
and all nine copying-machine rule constructors. Both term and atom sharing
sorts are included in the one template.

\lean{Copying/ContractionIdempotenceEquivalence}{10}{12}{ContractionIdempotence.template_robust}
\lean{Copying/ContractionIdempotenceEquivalence}{15}{17}{ContractionIdempotence.template_robust_converse}
Arbitrary closed recursive Idem fillings now have exact execution
equivalence and agree at every successful count. This includes many
simultaneous occurrences and occurrences inside boxes that are later opened
or copied. The comparison retains the literal two weakenings and contraction
on one side and the single weakening on the other. These graphs are not
identified by representation isomorphism: their primitive labels and edge
counts differ. The equality is instead the proved semantic conclusion about
successful machine observations.

\lean{Copying/ContractionIdempotenceEquivalence}{21}{25}{ContractionIdempotence.filling_exact_equiv}
\lean{Copying/ContractionIdempotenceEquivalence}{28}{32}{ContractionIdempotence.filling_termination_count}
Finally, each original Idem pair, including independent ordered carrier
renamings, can be substituted at any addressed position of any closed
observing context. Its two filled graphs have exact quantitative equivalence
and identical successful-transition counts. Since the pair has empty input
interface, this substantive context theorem is stated for addressed
fragment replacement, rather than incorrectly treating the pair itself as
a one-term-input open program. It supplies the full Idem dependency needed
by the box-weakening and later semantic weakening arguments; those further
arguments remain separate obligations.

\lean{Copying/ContractionIdempotenceEquivalence}{37}{42}{ContractionIdempotence.addressed_exact_equiv}
\lean{Copying/ContractionIdempotenceEquivalence}{46}{51}{ContractionIdempotence.addressed_termination_count}
\subsection{Concrete regressions for both sharing sorts and copied boxes}
Close a term output with the original unit constant and an atom output with
an original store followed by unit. The general label lemma for these
consumers follows by primitive and serial composition. Place the literal
Idem fragment followed by this consumer beside an executable constant-seven
root. The resulting observing graph has one term input and no outputs,
so it is an actual executable example for either wire sort.
\lean{Copying/ContractionIdempotenceExamples}{12}{14}{ContractionIdempotence.Examples.finish}
\lean{Copying/ContractionIdempotenceExamples}{16}{22}{ContractionIdempotence.Examples.finish_allLabels}
\lean{Copying/ContractionIdempotenceExamples}{26}{27}{ContractionIdempotence.Examples.program}
The literal pair is not recursively isomorphic: transporting its contraction
label to the one-edge right graph would identify contraction with weakening.
A genuine recursive witness places this pair in the stated observing context,
and primitive composition proves hole freedom. Its actual common context is
binding-free as a derived consequence of nullarity; the semantic theorem
still quantifies over unrestricted observations.
\lean{Copying/ContractionIdempotenceExamples}{31}{37}{ContractionIdempotence.Examples.left_not_iso}
\lean{Copying/ContractionIdempotenceExamples}{41}{47}{ContractionIdempotence.Examples.programWitness}
\lean{Copying/ContractionIdempotenceExamples}{49}{54}{ContractionIdempotence.Examples.program_holeFree}
\lean{Copying/ContractionIdempotenceExamples}{58}{60}{ContractionIdempotence.Examples.program_bindingFree}
\par\penalty0

For either fragment, the executable root performs exactly one original
search transition and is final in backtrack mode. The exact-cost recurrence
proves that every successful finite run has count one, including arbitrary
contextual rewrite presentations. Independently, applying the completed Idem
semantic theorem to the actual recursive witness proves exact unbounded
execution equivalence of the two closed examples.
\lean{Copying/ContractionIdempotenceExamples}{64}{67}{ContractionIdempotence.Examples.program_returns}
\lean{Copying/ContractionIdempotenceExamples}{71}{78}{ContractionIdempotence.Examples.program_exact_cost}
\lean{Copying/ContractionIdempotenceExamples}{82}{88}{ContractionIdempotence.Examples.program_exact_equiv}
A second recursive witness inserts each literal Idem pair into a genuine
lambda body with an unused bound output. Applying the established deep-copy
witness construction keeps the pair in both copied residual bodies, with the
original output distributors. The final theorem constructs the two actual
original copying steps for these lambdas, for both term and atom sharing.
It requires no evaluation of the body in advance. Together these regressions
check concrete executable contexts, the distinction from representation
isomorphism, and the copying/deep-box dependency of the generic proof.
\lean{Copying/ContractionIdempotenceExamples}{92}{93}{ContractionIdempotence.Examples.bodyWitness}
\lean{Copying/ContractionIdempotenceExamples}{97}{98}{ContractionIdempotence.Examples.copiedBodyWitness}
\lean{Copying/ContractionIdempotenceExamples}{102}{117}{ContractionIdempotence.Examples.body_copy_steps}